\chapter*{Conclusion générale}
\addcontentsline{toc}{chapter}{Conclusion générale}

Ce stage a permis de structurer une vision fonctionnelle cohérente d'un Smart WMS, à partir d'un benchmark des solutions du marché, de l'analyse des besoins d'entrepôt et de l'identification des cas d'utilisation. Le travail ne consistait pas à juxtaposer des technologies, mais à définir un socle WMS utilisable dans différents contextes logistiques et capable d'évoluer selon la maturité de l'organisation.

L'architecture proposée sépare clairement le coeur opérationnel, les capacités d'aide à la décision et les moyens techniques dépendants de l'infrastructure. Cette distinction 	extit{Core / Smart / Optional} est un résultat central du stage. Elle garantit que les opérations de réception, de stockage, d'inventaire, de préparation et d'expédition restent réalisables sans dépendre d'un dispositif avancé. Les fonctions de prévision, d'optimisation et de détection d'anomalies peuvent ensuite être activées lorsque les données disponibles justifient leur emploi. Les équipements tels que RFID, vision, OCR ou AMR/AGV restent des options intégrables sans modifier le périmètre métier fondamental.

La modélisation UML a permis de vérifier cette cohérence sous plusieurs angles. Les cas d'utilisation relient les acteurs aux fonctions attendues ; les activités rendent visibles les décisions et les exceptions ; les séquences précisent les échanges avec l'ERP, le TMS/YMS et les canaux terrain ; enfin, le modèle de domaine et l'organisation logique des composants établissent une base utile à la conception applicative future. Les diagrammes obtenus constituent ainsi des supports de communication entre équipes métier, logistique et informatique.

Les perspectives concernent la validation du modèle auprès d'utilisateurs opérationnels, la priorisation des cas d'utilisation selon la valeur métier et la disponibilité des données, puis la réalisation progressive d'un prototype. Cette étape devra s'appuyer sur des indicateurs mesurables : fiabilité des stocks, respect des délais, productivité, taux d'anomalies et consommation énergétique. L'approche proposée permet donc de faire évoluer l'entrepôt de manière maîtrisée, en plaçant la qualité des processus et des données avant la sophistication technologique.
